\section{Introduction}\label{sec:r9-intro}
An inspection service has to answer an engineering question beyond whether its latest pose estimate is accurate: does the available evidence support the requirement that the consumer currently asks it to satisfy? A camera can acquire an image with adequate sampling density while its target coverage is insufficient. A localization update can be useful for navigation yet arrive too late to geotag an earlier image. A report can meet its delivery requirement without satisfying the complete acquisition-and-association requirement. Reusing a statistical statement after these requirements change therefore requires an interpretation of the event being certified, its producer and the evidence lineage.

Engineering knowledge representations make requirements, domain quantities and change dependencies explicit \cite{petnga2016,wang2018change,ma2023inspection}. The decision considered here is whether to retain a calibrated event statement, re-evaluate archived calibration outputs, or refuse a request. This is an engineering evidence-management problem whose numerical producer is an anytime airframe-relative localization service. The application is not described as a digital twin, and a metadata match is not treated as evidence that an unknown deployment distribution remains unchanged.

The retained localization study already includes a rigorous control: a fully validated procedural system can reproduce the contract system's numerical and authorization results. Fixed2 also remains competitive on ordinary routes. Those observations motivate an additional question: what capability does explicit task-event knowledge add when the consumer's requirements or the producer's configuration change? Improved localization error is neither necessary nor sufficient to answer that question.

We develop a typed interface that compiles requirements into failure events, emits a pointwise event-containment witness, and records the producer dependencies needed for that particular requested event. A calibrated anchor can cover a family of weaker requirements without rescoring calibration records. A producer edit can retain such evidence only when the requested event's implemented functional dependencies remain invariant. Otherwise the system renews the evidence from applicable records or refuses it. The interval algebra, binomial inversion and dependency reasoning are established tools; the contribution is their executable engineering interface, its explicit transfer conditions and its evaluation under requirement evolution.

Two prospective studies add image acquisition, capture-time spatial association and report delivery to the retained geometric producer. They use actual synthetic rasters and pixel processing, independently generated inputs, unchanged numerical predictors, and a predeclared wider-ROI intervention. A complete procedural semantic implementation and a cached direct-renewal baseline receive the same information and confidence budget. Their equality, superior support or lower runtime, when observed, remains part of the result. The study distinguishes requirement reuse, producer-change applicability and numerical scheduling; it does not claim an exclusive advantage of graph storage or a physical robot risk guarantee.

The paper first specifies the engineering quantities and retained executor, then presents the typed event interface and its conditions for evidence reuse. The prospective design and results evaluate service support, refusal, dependency propagation, actual record re-evaluation and computational trade-offs. Historical source and route results are retained separately to establish what the localization producer can and cannot provide.
